\section{Chapter~\ref{sec:code}. Morse Code}

The chapter's limiting estimates are information theory and arithmetic; what has been measured is where noise arises in the channel, how fast the brain works, and what a spike costs, while the question of which code the cortex actually uses remains open.

{\footnotesize
\setlength{\tabcolsep}{3pt}\renewcommand{\arraystretch}{1.15}
\begin{longtable}{>{\raggedright}p{0.5cm}>{\raggedright}p{4.0cm}>{\raggedright}p{5.6cm}>{\raggedright}p{2.3cm}>{\raggedright\arraybackslash}p{1.7cm}}
\toprule
No. & Claim & Experiment and what was measured & Source & Status \\
\midrule\endfirsthead
\toprule
No. & Claim & Experiment and what was measured & Source & Status \\
\midrule\endhead
1 & Rates of cortical \nt{GLUT} neurons range from fractions to tens of spikes per second, and most of the time neurons work near the lower end & Cell-attached recording (neuron selection independent of its activity) in the auditory cortex of awake rats: at any moment fewer than $5\,\%$ of neurons fire above $20$ spikes per second; some neurons have a background rate below $0.01$ per second; the rate distribution is lognormal & \cite{Hromadka2008} & measured \\
2 & The capacity of a spike channel is $R_{\max}\approx r\log_2\frac{e}{r\Delta t}$~\eqref{eq:spikeentropy}, and almost all of it lies in spike timing (Proposition~\ref{prop:capacity}) & Entropy of a binary string of bins of length $\Delta t$. The chapter's numbers: $8$ bits per spike at $r=10$ per second and $\Delta t=1\ms$, about $5$ bits at $\Delta t=10\ms$, fewer than $2$ bits per second for a spike counter & \cite{MacKay1952} & theory \\
3 & Sensory neurons transmit one to several bits per spike, in the best cases up to half the limit & Stimulus reconstruction from the spikes of sensory neurons whose stimulus is known; a summary of measurements in the book & \cite{Rieke1997} & measured \\
4 & The spike generator is precise; precise timing is set by fast transients of the input & Rat cortical neurons in slices: for a repeated current with fast fluctuations, spikes repeat with a precision better than $1\ms$; for a constant current the times drift & \cite{Mainen1995} & measured \\
5 & The synapse is unreliable: more than half of the spikes do not reach the recipient because of random release & Minimal stimulation of inputs to hippocampal pyramidal neurons (slices): more than half of the spikes evoke no response. Fluctuations of the stimulation threshold, conduction failures at axon branch points, and low experimental temperature were ruled out; what remains is probabilistic release & \cite{AllenStevens1994} & measured \\
6 & The spike train in living cortex looks random: intervals are scattered as in a Poisson stream, and the variance of the spike count is close to the mean; part of the ``noise'' is a message about the animal's movements & Recordings in visual areas V1 and MT of awake monkeys: coefficient of variation of the intervals about $1$, ratio of spike-count variance to mean as for a random process; a model summing independent small inputs gives much less variability. In mouse visual cortex a considerable part of the variability is predicted from video of the animal's own movements & \cite{Softky1993,Stringer2019} & measured \\
7 & The rate code is slow: error $1/\sqrt{rT}$~\eqref{eq:rateerror}, while the brain recognizes a picture in $\sim150\ms$ & The formula is Poisson statistics, the chapter's derivation. Experiment: people decided whether an unfamiliar photograph flashed for $20\ms$ contained an animal; the brain potentials for ``yes'' and ``no'' diverge after about $150\ms$. Splitting this time into ten stages of $10$--$20\ms$ is the chapter's estimate, not a measurement & \cite{Thorpe1996} & measured \\
8 & Averaging over a group is limited by correlation: the error falls by no more than a factor of $1/\sqrt c$~\eqref{eq:correlated} (Proposition~\ref{prop:pooling}) & Pairs of monkey MT neurons recorded simultaneously: the spike-count correlation coefficient averages $0.12$, and theoretical analysis shows that even such a correlation limits the gain from a group. Theory: only the part of the correlation that looks like the signal sets the limit. A model of the opposing position: the rate of a group of $50$--$100$ neurons is read within one interspike interval, $10$--$50\ms$, and the timing of individual spikes is not used in the cortex & \cite{Zohary1994,MorenoBote2014,ShadlenNewsome1998} & measured \\
9 & A number can be written in the first-spike time~\eqref{eq:latency}, in the spike order of a group, or in the phase of a local rhythm & Salamander retina: the spatial structure of a briefly flashed image is recorded in the relative latencies of the first spikes of the output neurons, almost independently of contrast. Rat hippocampal place cells: while passing through ``their'' place, spikes shift earlier and earlier in the theta cycle ($7$--$12$\,Hz), a shift of $100^\circ$ to $355^\circ$. How much the cortex uses spike timing is an open question (row~8) & \cite{Gollisch2008,OKeefe1993} & measured \\
10 & Which code gets read is decided by the recipient's time constant~\eqref{eq:reader} (Proposition~\ref{prop:reader}); in active cortex neurons are closer to coincidence detectors & Formula~\eqref{eq:reader} is the chapter's derivation. Review of in vivo recordings: in active cortex the membrane conductance is several times higher than at rest, and the time constant correspondingly shorter. That the cortex switches codes in precisely this way has not been shown directly & \cite{Destexhe2003} & hypothesis \\
11 & A depressing synapse transmits the change in rate, essentially $\Delta r/r$~\eqref{eq:depression} (Fig.~\ref{fig:depression}) & Paired recordings of cortical pyramidal neurons: the rate of depression is set by the release probability; with slow depression the response reflects rate, with fast depression coincidences. A model based on measured depression in rat cortex: equal relative rate changes on fast and slow inputs give equal responses, that is, the synapse transmits $\Delta r/r$. The numbers $1.7\to2.2$, $6.7$, and $90\ms$ are the chapter's calculation & \cite{Tsodyks1997,Abbott1997depression} & measured \\
12 & A burst is a ``dash'': its probability of being lost, $(1-p)^k$~\eqref{eq:burst}, is small, and facilitation lowers it further & Review: at unreliable synapses single spikes are often lost, while bursts are transmitted reliably thanks to facilitation; a burst induces long-term synaptic changes. That the brain reads a burst as a separate symbol is a hypothesis & \cite{Lisman1997,AllenStevens1994} & hypothesis \\
13 & Fast \nt{GABA} neurons set the rhythm and ``punctuation''; one more class inhibits the inhibitors and opens the gates (Proposition~\ref{prop:punctuation}) & Review of the properties of classes of cortical \nt{GABA} neurons. Optogenetics: rhythmic excitation of fast \nt{GABA} neurons evokes the gamma rhythm, and the response to a stimulus depends on the phase of the cycle. In mouse visual cortex PV neurons inhibit one another, SST neurons inhibit all the other classes, and VIP neurons inhibit mainly SST. Excitation of VIP neurons in awake mice releases \nt{GLUT} neurons from inhibition; they are switched on by a reinforcement signal. The division of roles as a general rule is a hypothesis & \cite{Tremblay2016,Cardin2009,Pfeffer2013,Pi2013} & hypothesis \\
14 & Energy limits the mean rate to about one spike per second~\eqref{eq:energy} & Energy budget of rodent gray matter from anatomical and physiological data: spikes and synaptic currents take $47\,\%$ and $34\,\%$ of the signaling energy; no more than $\sim15\,\%$ of neurons can be active at once. For the human cortex: no more than $\sim1\,\%$ of neurons can be strongly active at once. The figures of $\sim10^9$ ATP per spike and $\sim10$\,W for signaling are the chapter's estimate based on these calculations & \cite{Attwell2001,Lennie2003} & theory \\
15 & The code is sparse and tuned for the greatest number of bits per joule; the cortex trades precision for energy (Proposition~\ref{prop:economy}) & The efficient-coding hypothesis. Support: in food-restricted mice, AMPA receptor conductance and synaptic ATP use in the visual cortex fall ($-29\,\%$), the firing rate is preserved, and orientation tuning broadens by $32\,\%$ & \cite{Barlow1961,Padamsey2022} & hypothesis \\
\bottomrule
\end{longtable}}
